\section{Refleksion over vejledning}
Gruppen har oplevet, at vejledningen har fungeret godt i den del af forløbet, hvor der stadig har været planlagt undervisning og faste vejledningstimer. I denne periode har det været muligt at få støtte og sparring, hvilket har bidraget positivt til projektets fremdrift.
\\\\
Samtidig har gruppen oplevet, at det kan være begrænsende kun at have en enkelt vejleder tilknyttet projektet, da vejlederen ikke nødvendigvis har haft faglig indsigt inden for alle relevante områder. Dette har i nogle tilfælde gjort det vanskeligt at få helt konkret og brugbar feedback på mere specialiserede problemstillinger, særligt i forhold til de Formaila aspekter af projektet.
\\\\
Efter den planlagte undervisnings- og vejledningsperiode har gruppen desuden erfaret, at det kræver en mere proaktiv tilgang selv at tage kontakt til vejlederne ved opståede problemer eller tvivlsspørgsmål. Gruppen har i den forbindelse lært, at det er vigtigt aktivt at opsøge vejledning løbende, frem for at afvente faste møder.
\\\\
Fremadrettet vil gruppen derfor fokusere på at være mere opsøgende og struktureret i brugen af vejledere, blandt andet ved at forberede konkrete spørgsmål på forhånd og tage kontakt tidligere i processen, når udfordringer opstår. Dette forventes at kunne forbedre udbyttet af vejledningen og bidrage til en mere effektiv projektudvikling.